\documentclass[11pt, a4paper]{report}

% ========== Common Packages ==========
\usepackage{fontspec}
\usepackage{kotex}
\usepackage{amsmath, amssymb, amsthm}
\usepackage{mathtools}
\usepackage{bm}
\usepackage{graphicx}
\usepackage{booktabs}
\usepackage{array}
\usepackage{tabularx}
\usepackage{longtable}
\usepackage{hyperref}
\usepackage{cleveref}
\usepackage{geometry}
\usepackage{fancyhdr}
\usepackage{titlesec}
\usepackage{enumitem}
\usepackage{xcolor}
\usepackage{tcolorbox}
\usepackage{algorithm}
\usepackage{algorithmic}
\usepackage{tikz}
\usetikzlibrary{arrows.meta, positioning, calc, decorations.pathreplacing, shapes.geometric, fit, patterns, angles, quotes, trees}
\usepackage{pgfplots}
\pgfplotsset{compat=1.18}
\usepackage{caption}
\usepackage{subcaption}
\usepackage{float}

\geometry{margin=2.5cm, top=3cm, bottom=3cm}

\usepackage{multirow}
\usepackage{siunitx}
\usepackage{listings}
% ========== Colors ==========
% Common
\definecolor{defblue}{RGB}{30, 80, 150}
\definecolor{thmgreen}{RGB}{20, 110, 60}
\definecolor{noteorange}{RGB}{200, 100, 0}
\definecolor{lightblue}{RGB}{230, 240, 255}
\definecolor{lightgreen}{RGB}{230, 255, 235}
\definecolor{lightorange}{RGB}{255, 245, 230}
\definecolor{lightgray}{RGB}{245, 245, 245}
% Extended
\definecolor{lightred}{RGB}{255, 235, 235}
\definecolor{darkred}{RGB}{180, 30, 30}
\definecolor{biogreen}{RGB}{10, 130, 80}
\definecolor{cvpurple}{RGB}{100, 40, 150}
\definecolor{injuryred}{RGB}{200, 40, 40}
\definecolor{codegray}{RGB}{240, 240, 245}
\definecolor{pipeblue}{RGB}{25, 100, 180}
\definecolor{constraintred}{RGB}{200, 50, 50}
\definecolor{termblack}{RGB}{30, 30, 30}
\definecolor{goldcolor}{RGB}{180, 140, 20}

% ========== tcolorbox environments ==========
% --- Common (all parts) ---
\newtcolorbox{keyidea}[1][]{
  colback=lightblue,
  colframe=defblue,
  fonttitle=\bfseries,
  title={Key Idea},
  #1
}

\newtcolorbox{importantnote}[1][]{
  colback=lightorange,
  colframe=noteorange,
  fonttitle=\bfseries,
  title={Important},
  #1
}

\newtcolorbox{summary}[1][]{
  colback=lightgreen,
  colframe=thmgreen,
  fonttitle=\bfseries,
  title={Summary},
  #1
}

% --- Warning/Error ---
\newtcolorbox{warningbox}[1][]{
  colback=lightred,
  colframe=darkred,
  fonttitle=\bfseries,
  title={Warning},
  #1
}

% --- Domain: Biomechanics & Clinical ---
\newtcolorbox{clinicalnote}[1][]{
  colback=lightred,
  colframe=darkred,
  fonttitle=\bfseries,
  title={Clinical Note},
  #1
}

\newtcolorbox{biomechbox}[1][]{
  colback=lightgreen,
  colframe=biogreen,
  fonttitle=\bfseries,
  title={Biomechanics},
  #1
}

\newtcolorbox{injurybox}[1][]{
  colback={injuryred!8},
  colframe=injuryred,
  fonttitle=\bfseries,
  title={Injury Risk},
  #1
}

% --- Domain: Computer Vision ---
\newtcolorbox{cvbox}[1][]{
  colback={cvpurple!5},
  colframe=cvpurple,
  fonttitle=\bfseries,
  title={Computer Vision},
  #1
}

% --- Pipeline & Setup ---
\newtcolorbox{setupbox}[1][]{
  colback=lightgray,
  colframe=gray!60,
  fonttitle=\bfseries,
  title={Setup},
  #1
}

\newtcolorbox{pipelinebox}[1][]{
  colback=pipeblue!5,
  colframe=pipeblue,
  fonttitle=\bfseries,
  title={Pipeline Step},
  #1
}

\newtcolorbox{codebox}[1][]{
  colback=codegray,
  colframe=gray!70,
  fonttitle=\bfseries\ttfamily,
  title={Code},
  #1
}

\newtcolorbox{termbox}[1][]{
  colback=termblack!5,
  colframe=termblack!60,
  fonttitle=\bfseries\ttfamily,
  title={Terminal},
  #1
}

% --- Research ---
\newtcolorbox{hypothesis}[1][]{
  colback=goldcolor!8,
  colframe=goldcolor,
  fonttitle=\bfseries,
  title={Hypothesis},
  #1
}

\newtcolorbox{resultbox}[1][]{
  colback=thmgreen!8,
  colframe=thmgreen,
  fonttitle=\bfseries,
  title={Expected Result},
  #1
}

% --- Constraints & Solutions ---
\newtcolorbox{constraintbox}[1][]{
  colback=constraintred!8,
  colframe=constraintred,
  fonttitle=\bfseries,
  title={Constraint},
  #1
}

\newtcolorbox{solutionbox}[1][]{
  colback=thmgreen!8,
  colframe=thmgreen,
  fonttitle=\bfseries,
  title={Solution},
  #1
}

\newtcolorbox{databox}[1][]{
  colback=pipeblue!5,
  colframe=pipeblue,
  fonttitle=\bfseries,
  title={Data Spec},
  #1
}

% --- Progress/Status ---
\newtcolorbox{successbox}[1][]{
  colback=lightgreen,
  colframe=thmgreen,
  fonttitle=\bfseries,
  title={Success},
  #1
}

\newtcolorbox{nextbox}[1][]{
  colback=lightorange,
  colframe=noteorange,
  fonttitle=\bfseries,
  title={Next Step},
  #1
}

% --- Fitting guide ---
\newtcolorbox{failbox}[1][]{
  colback=lightred,
  colframe=darkred,
  fonttitle=\bfseries,
  title={Failure Pattern},
  #1
}

\newtcolorbox{fixbox}[1][]{
  colback=lightgreen,
  colframe=thmgreen,
  fonttitle=\bfseries,
  title={Fix},
  #1
}

\newtcolorbox{lessonbox}[1][]{
  colback=lightorange,
  colframe=noteorange,
  fonttitle=\bfseries,
  title={Lesson Learned},
  #1
}

% --- Specification ---
\newtcolorbox{specbox}[1][]{
  colback=cvpurple!5,
  colframe=cvpurple,
  fonttitle=\bfseries,
  title={Specification},
  #1
}

\newtcolorbox{checklistbox}[1][]{
  colback=lightgreen,
  colframe=biogreen,
  fonttitle=\bfseries,
  title={Checklist},
  #1
}

% ========== Custom math commands ==========
% Vectors (lowercase)
\newcommand{\vx}{\mathbf{x}}
\newcommand{\vd}{\mathbf{d}}
\newcommand{\vo}{\mathbf{o}}
\newcommand{\vc}{\mathbf{c}}
\newcommand{\vr}{\mathbf{r}}
\newcommand{\vp}{\mathbf{p}}
\newcommand{\vv}{\mathbf{v}}
\newcommand{\vn}{\mathbf{n}}
\newcommand{\vu}{\mathbf{u}}
\newcommand{\vt}{\mathbf{t}}
\newcommand{\ve}{\mathbf{e}}
\newcommand{\vq}{\mathbf{q}}
% Vectors (uppercase)
\newcommand{\vF}{\mathbf{F}}
\newcommand{\vM}{\mathbf{M}}
\newcommand{\va}{\mathbf{a}}
\newcommand{\vg}{\mathbf{g}}
\newcommand{\vX}{\mathbf{X}}
% Bold Greek
\newcommand{\vmu}{\bm{\mu}}
\newcommand{\vbeta}{\bm{\beta}}
\newcommand{\vtheta}{\bm{\theta}}
\newcommand{\vpsi}{\bm{\psi}}
% Matrices
\newcommand{\mSigma}{\bm{\Sigma}}
\newcommand{\mR}{\mathbf{R}}
\newcommand{\mS}{\mathbf{S}}
\newcommand{\mW}{\mathbf{W}}
\newcommand{\mJ}{\mathbf{J}}
\newcommand{\mG}{\mathbf{G}}
\newcommand{\mT}{\mathbf{T}}
\newcommand{\mI}{\mathbf{I}}
\newcommand{\mB}{\mathbf{B}}
\newcommand{\mK}{\mathbf{K}}
\newcommand{\mP}{\mathbf{P}}
\newcommand{\mF}{\mathbf{F}}
\newcommand{\mE}{\mathbf{E}}
\newcommand{\mA}{\mathbf{A}}
\newcommand{\mH}{\mathbf{H}}
% Sets and groups
\newcommand{\RR}{\mathbb{R}}
\newcommand{\SO}{\mathrm{SO}}
% Units
\newcommand{\degps}{\,\text{deg/s}}
\newcommand{\Nm}{\ensuremath{\,\text{N}\cdot\text{m}}}
\newcommand{\BW}{\,\text{BW}}

% ========== Header/Footer ==========
\pagestyle{fancy}
\fancyhf{}
\fancyhead[L]{\leftmark}
\fancyhead[R]{\thepage}
\renewcommand{\headrulewidth}{0.4pt}

% ========== Title formatting ==========
\titleformat{\chapter}[display]
  {\normalfont\huge\bfseries\color{defblue}}
  {\chaptertitlename\ \thechapter}{20pt}{\Huge}

\titleformat{\section}
  {\normalfont\Large\bfseries\color{defblue!80!black}}
  {\thesection}{1em}{}

\titleformat{\subsection}
  {\normalfont\large\bfseries\color{defblue!60!black}}
  {\thesubsection}{1em}{}


\definecolor{codegray}{RGB}{240, 240, 245}
\lstset{
  basicstyle=\ttfamily\small,
  keywordstyle=\color{defblue}\bfseries,
  commentstyle=\color{thmgreen}\itshape,
  stringstyle=\color{noteorange},
  backgroundcolor=\color{codegray},
  frame=single,
  rulecolor=\color{gray!40},
  breaklines=true,
  showstringspaces=false,
  numbers=left,
  numberstyle=\tiny\color{gray},
  tabsize=4,
  captionpos=b,
  language=Python
}

\begin{document}

\begin{titlepage}
\centering
\vspace*{1.5cm}
{\Large\color{gray} Integration Guide}
\vspace{0.5cm}

{\Huge\bfseries\color{defblue} VGGT + EasyMocap\\[0.3cm]
통합 파이프라인}

\vspace{1cm}

{\LARGE\color{defblue!70!black} VGGT 카메라 캘리브레이션 결과를\\
EasyMocap SMPL 피팅에 활용하는\\
방법론, 이론, 구현}

\vspace{1.5cm}

\begin{tikzpicture}[
  block/.style={rectangle, rounded corners=4pt, draw=defblue, fill=lightblue,
    minimum width=2.2cm, minimum height=0.7cm, font=\small\bfseries, text=defblue!80!black, align=center},
  done/.style={block, draw=thmgreen, fill=thmgreen!15, text=thmgreen!80!black},
  curr/.style={block, draw=noteorange, fill=noteorange!15, text=noteorange!80!black},
  arrow/.style={-Stealth, thick, defblue!60}
]
\node[done] (vggt) at (0,0) {VGGT\\K, R, t};
\node[curr] (conv) at (3.5,0) {Format\\Converter};
\node[curr] (emc) at (7,0) {EasyMocap\\SMPL Fitting};
\node[block] (gart) at (10.5,0) {GART\\3DGS};

\draw[arrow, thmgreen] (vggt) -- (conv);
\draw[arrow, noteorange] (conv) -- (emc);
\draw[arrow, gray!50] (emc) -- (gart);

\node[font=\tiny\color{thmgreen}] at (0, -0.6) {완료};
\node[font=\tiny\color{noteorange}] at (3.5, -0.6) {본 문서};
\node[font=\tiny\color{noteorange}] at (7, -0.6) {실행 대기};
\node[font=\tiny\color{gray}] at (10.5, -0.6) {다음 단계};
\end{tikzpicture}

\vspace{2cm}
{\large 2026년 4월 12일}

\vspace{\fill}
{\small\color{gray} VGGT (CVPR 2025)의 카메라 파라미터를 EasyMocap의 intri.yml / extri.yml 포맷으로\\
변환하여, 기존에 좌표계 불일치로 실패했던 SMPL 피팅을 재시도합니다.}
\end{titlepage}

\tableofcontents
\newpage


% ====================================================================
\chapter{문제 정의와 전략}
% ====================================================================

\section{이전 실패 원인 요약}

\begin{warningbox}
이전 EasyMocap SMPL 피팅이 실패한 핵심 원인:
\begin{enumerate}[leftmargin=*]
  \item DA3의 K가 280$\times$280에서 추정되어, 원본 해상도로의 역변환 과정에서 $s_x \neq s_y$ 발생
  \item 역변환된 K와 OpenPose 2D 키포인트의 좌표계 불일치
  \item 결과적으로 삼각측량된 3D 점과 SMPL 포즈가 다른 좌표 공간에 존재
  \item EasyMocap 출력의 Y/Z 축 교환 현상 (Th의 좌표가 기대값과 불일치)
\end{enumerate}
\end{warningbox}

\section{새 전략: VGGT $\rightarrow$ EasyMocap}

\begin{keyidea}
VGGT의 K는 원본 해상도 좌표계에서 직접 매칭되므로, OpenPose 2D 키포인트와 좌표계가 자동으로 일치한다. 이 K와 R, t를 EasyMocap의 포맷으로 변환하면, 좌표계 불일치 문제가 구조적으로 해결된다.
\end{keyidea}

\begin{figure}[H]
\centering
\begin{tikzpicture}[
  block/.style={rectangle, rounded corners=4pt, draw=defblue, fill=lightblue,
    minimum width=3.5cm, minimum height=1cm, font=\small, text=defblue!80!black, align=center},
  file/.style={rectangle, rounded corners=2pt, draw=gray!60, fill=lightgray,
    minimum width=2.5cm, minimum height=0.6cm, font=\small\ttfamily, align=center},
  arrow/.style={-Stealth, thick}
]
% VGGT
\node[block] (vggt) at (0, 3) {\textbf{VGGT Inference}\\7장 이미지 $\rightarrow$ K, R, t};
\node[file] (vjson) at (0, 1.5) {vggt\_result.json};

% Converter
\node[block] (conv) at (0, 0) {\textbf{Format Converter}\\VGGT $\rightarrow$ OpenCV YAML};
\node[file] (intri) at (-2.5, -1.5) {intri.yml};
\node[file] (extri) at (2.5, -1.5) {extri.yml};

% EasyMocap
\node[block] (emc) at (0, -3) {\textbf{EasyMocap}\\triangulate + fitSMPL};
\node[file] (smpl) at (0, -4.5) {smpl/000000.json};

% OpenPose
\node[file] (op) at (6, -3) {annots/\{1..7\}/\\*.json};

\draw[arrow, defblue] (vggt) -- (vjson);
\draw[arrow, defblue] (vjson) -- (conv);
\draw[arrow, defblue] (conv) -- (intri);
\draw[arrow, defblue] (conv) -- (extri);
\draw[arrow, noteorange] (intri) -- (emc);
\draw[arrow, noteorange] (extri) -- (emc);
\draw[arrow, gray] (op) -- (emc);
\draw[arrow, thmgreen] (emc) -- (smpl);
\end{tikzpicture}
\caption{VGGT $\rightarrow$ EasyMocap 통합 파이프라인}
\end{figure}


% ====================================================================
\chapter{이론: 좌표계와 변환}
% ====================================================================

\section{VGGT의 좌표 규약}

VGGT는 \textbf{OpenCV 규약}을 사용한다:
\begin{itemize}[leftmargin=*]
  \item $x$: 오른쪽, $y$: 아래쪽, $z$: 앞쪽 (카메라 시선 방향)
  \item Extrinsic $[\mR \mid \vt]$는 \textbf{world-to-camera} 변환
  \item $\vp_\text{cam} = \mR \cdot \vp_\text{world} + \vt$
\end{itemize}

카메라의 월드 좌표 위치는:
\begin{equation}
\vp_\text{cam\_center} = -\mR^\top \vt
\end{equation}

\section{EasyMocap의 좌표 규약}

EasyMocap도 \textbf{OpenCV 규약}을 사용한다:
\begin{itemize}[leftmargin=*]
  \item \texttt{extri.yml}의 \texttt{Rot\_X}: $3 \times 3$ 회전 행렬 (world-to-camera)
  \item \texttt{extri.yml}의 \texttt{T\_X}: $3 \times 1$ 평행이동 벡터
  \item \texttt{extri.yml}의 \texttt{R\_X}: Rodrigues 벡터 ($3 \times 1$, 축-각도 표현)
  \item \texttt{intri.yml}의 \texttt{K\_X}: $3 \times 3$ 내부 행렬
  \item \texttt{intri.yml}의 \texttt{dist\_X}: $1 \times 5$ 왜곡 계수 ($k_1, k_2, p_1, p_2, k_3$)
\end{itemize}

\begin{summary}
\textbf{VGGT와 EasyMocap은 같은 좌표 규약 (OpenCV)}을 사용한다.
따라서 좌표 변환이 필요 없고, 포맷만 변환하면 된다.
\end{summary}


\section{투영 모델}

두 시스템 모두 동일한 핀홀 투영 모델을 사용:
\begin{equation}
\lambda \begin{pmatrix} u \\ v \\ 1 \end{pmatrix}
= \underbrace{\begin{pmatrix} f_x & 0 & c_x \\ 0 & f_y & c_y \\ 0 & 0 & 1 \end{pmatrix}}_{\mK}
\underbrace{\begin{pmatrix} \mR & \vt \end{pmatrix}}_{[\mR|\vt]_\text{3$\times$4}}
\begin{pmatrix} X \\ Y \\ Z \\ 1 \end{pmatrix}
\end{equation}

여기서 $(u, v)$는 OpenPose가 검출한 2D 키포인트 좌표 (원본 해상도 기준).
VGGT의 K가 원본 해상도 기준이므로, $(u, v)$와 K가 같은 좌표계에 있다.


\section{EasyMocap의 삼각측량}

EasyMocap의 \texttt{SimpleTriangulate}는 다음을 수행:

\begin{enumerate}[leftmargin=*]
  \item 각 카메라 $i$에서 투영 행렬 구성: $\mP_i = \mK_i [\mR_i \mid \vt_i]$
  \item 2D 키포인트 $(u_j^i, v_j^i)$와 $\mP_i$로 DLT 삼각측량
  \item Iterative refinement: confidence 가중 + 재투영 오차 기반 outlier 제거
  \item 결과: 3D 키포인트 $\vp_j \in \RR^3$ (월드 좌표)
\end{enumerate}

\begin{importantnote}
삼각측량의 정확도는 전적으로 \textbf{K와 2D 키포인트의 좌표계 일치}에 의존한다.
이전에 DA3의 K가 틀려서 삼각측량이 부정확했고, 그 부정확한 3D 점으로 SMPL을 피팅하니 실패한 것이다.
VGGT의 K를 사용하면 이 문제가 해결된다.
\end{importantnote}


\section{EasyMocap의 SMPL 피팅}

삼각측량된 3D 키포인트로부터 SMPL 파라미터를 추정:

\begin{equation}
\min_{\vtheta, \vbeta, \mR_h, \vt_h} \underbrace{\sum_{j} w_j \| J_j(\vtheta, \vbeta) - \vp_j \|^2}_{\text{3D joint loss}}
+ \underbrace{\lambda_\theta \|\vtheta\|^2}_{\text{pose prior}}
+ \underbrace{\lambda_\beta \|\vbeta\|^2}_{\text{shape prior}}
\end{equation}

여기서:
\begin{itemize}[leftmargin=*]
  \item $J_j(\vtheta, \vbeta)$: SMPL의 forward kinematics로 계산된 관절 위치
  \item $\vp_j$: 삼각측량된 3D 키포인트 (25개, Body25 포맷)
  \item $\mR_h, \vt_h$: global rotation (Rodrigues), global translation
  \item $\vtheta \in \RR^{69}$: body pose (23 joints $\times$ 3 axis-angle)
  \item $\vbeta \in \RR^{10}$: body shape
\end{itemize}

EasyMocap의 추가 단계:
\begin{enumerate}[leftmargin=*]
  \item \textbf{Init}: 삼각측량된 pelvis 위치로 $\vt_h$ 초기화
  \item \textbf{Shape optimization}: $\vbeta$ 먼저 최적화 (골격 비율)
  \item \textbf{Pose optimization}: $\vtheta$ 최적화
  \item \textbf{Joint optimization}: 모든 파라미터 동시 최적화
  \item (Optional) \textbf{2D reprojection refinement}: 각 카메라의 2D 키포인트에 대한 재투영 오차도 추가
\end{enumerate}


% ====================================================================
\chapter{개발 방향}
% ====================================================================

\section{Phase 1: VGGT $\rightarrow$ EasyMocap 포맷 변환}

\begin{setupbox}[title={변환 요구사항}]
\begin{enumerate}[leftmargin=*]
  \item VGGT JSON $\rightarrow$ \texttt{intri.yml}: 카메라별 K, dist (왜곡 = 0)
  \item VGGT JSON $\rightarrow$ \texttt{extri.yml}: 카메라별 R (Rodrigues), Rot (회전행렬), T (평행이동)
  \item 카메라 이름 매핑: cam1$\rightarrow$``1'', cam2$\rightarrow$``2'', ...
\end{enumerate}
\end{setupbox}

\texttt{intri.yml} 포맷:
\begin{lstlisting}[language={}]
%YAML:1.0
---
names: ["1", "2", "3", "4", "5", "6", "7"]
K_1: !!opencv-matrix
  rows: 3
  cols: 3
  dt: d
  data: [fx, 0, cx, 0, fy, cy, 0, 0, 1]
dist_1: !!opencv-matrix
  rows: 1
  cols: 5
  dt: d
  data: [0, 0, 0, 0, 0]
\end{lstlisting}

\texttt{extri.yml} 포맷:
\begin{lstlisting}[language={}]
R_1: !!opencv-matrix     # Rodrigues vector (3x1)
  rows: 3
  cols: 1
  dt: d
  data: [rx, ry, rz]
Rot_1: !!opencv-matrix   # Rotation matrix (3x3)
  rows: 3
  cols: 3
  dt: d
  data: [r00, r01, r02, r10, r11, r12, r20, r21, r22]
T_1: !!opencv-matrix     # Translation (3x1)
  rows: 3
  cols: 1
  dt: d
  data: [tx, ty, tz]
\end{lstlisting}


\section{Phase 2: EasyMocap 실행}

\begin{pipelinebox}[title={실행 명령}]
\begin{lstlisting}[language=bash]
cd ~/EasyMocap
python3 apps/demo/mv1p.py \
  ~/easymocap_data \
  --out ~/easymocap_data/output/vggt_smpl \
  --cfg config/exp/mv1p/detect_triangulate_fitSMPL.yml \
  --vis_det --vis_repro
\end{lstlisting}
\end{pipelinebox}


\section{Phase 3: 검증}

\begin{enumerate}[leftmargin=*]
  \item \textbf{삼각측량 검증}: 재투영 오차 (reprojection error) 계산
    \begin{equation}
    e_\text{reproj} = \frac{1}{N \cdot C} \sum_{j=1}^{N} \sum_{i=1}^{C} \| \pi(\mK_i, \mR_i, \vt_i, \vp_j) - (u_j^i, v_j^i) \|
    \end{equation}
    목표: $e_\text{reproj} < 10$ px

  \item \textbf{SMPL 피팅 검증}: MPJPE 계산
    \begin{equation}
    \text{MPJPE} = \frac{1}{J} \sum_{j=1}^{J} \| J_j^\text{SMPL} - \vp_j^\text{triangulated} \|
    \end{equation}
    목표: MPJPE $< 50$ mm (이전 실패: 219.6 mm)

  \item \textbf{시각적 검증}: SMPL 메시를 각 카메라 뷰에 재투영하여 오버레이
\end{enumerate}


\section{Phase 4: 스케일 보정 (필요시)}

\begin{warningbox}
VGGT의 translation은 정규화된 좌표 ($\|t\| \sim 1$)이므로, SMPL의 크기가 현실과 다를 수 있다. SMPL의 pelvis 높이가 실제 피사체와 일치하도록 global scale factor를 곱해야 할 수 있다.

\begin{equation}
s = \frac{h_\text{subject\_cm} / 100}{\|J_\text{head} - J_\text{ankle}\|_\text{SMPL}}
\end{equation}
\end{warningbox}


% ====================================================================
\chapter{구현 코드}
% ====================================================================

\section{VGGT $\rightarrow$ EasyMocap 변환 스크립트}

\begin{lstlisting}[caption={vggt\_to\_easymocap.py --- VGGT 결과를 EasyMocap intri.yml/extri.yml로 변환}, label=lst:converter]
#!/usr/bin/env python3
"""
Convert VGGT calibration JSON to EasyMocap intri.yml / extri.yml.

Usage:
  python vggt_to_easymocap.py \
    --input vggt_calibration_result.json \
    --output_dir ~/easymocap_data
"""

import json
import argparse
import numpy as np
import cv2
from pathlib import Path


def write_opencv_matrix(f, name, mat, rows, cols):
    """Write a matrix in OpenCV YAML format."""
    f.write(f"{name}: !!opencv-matrix\n")
    f.write(f"  rows: {rows}\n")
    f.write(f"  cols: {cols}\n")
    f.write(f"  dt: d\n")
    flat = mat.flatten().tolist()
    data_str = ", ".join(f"{v:.10f}" for v in flat)
    f.write(f"  data: [{data_str}]\n")


def rotation_matrix_to_rodrigues(R):
    """Convert 3x3 rotation matrix to Rodrigues vector."""
    rvec, _ = cv2.Rodrigues(np.array(R, dtype=np.float64))
    return rvec.flatten()


def main():
    parser = argparse.ArgumentParser()
    parser.add_argument("--input", required=True,
                        help="VGGT calibration JSON")
    parser.add_argument("--output_dir", required=True,
                        help="EasyMocap data directory")
    args = parser.parse_args()

    with open(args.input) as f:
        vggt = json.load(f)

    out = Path(args.output_dir)
    cam_names = sorted(vggt.keys(),
                       key=lambda x: int(x.replace("cam", "")))
    names = [c.replace("cam", "") for c in cam_names]

    # ---- Write intri.yml ----
    intri_path = out / "intri.yml"
    with open(intri_path, "w") as f:
        f.write("%YAML:1.0\n---\n")
        name_list = ", ".join(f'"{n}"' for n in names)
        f.write(f"names:\n")
        for n in names:
            f.write(f'  - "{n}"\n')

        for cam_key, name in zip(cam_names, names):
            cam = vggt[cam_key]
            K = np.array(cam["K"], dtype=np.float64)
            write_opencv_matrix(f, f"K_{name}", K, 3, 3)
            dist = np.zeros(5)
            write_opencv_matrix(f, f"dist_{name}", dist, 1, 5)

    print(f"Written: {intri_path}")

    # ---- Write extri.yml ----
    extri_path = out / "extri.yml"
    with open(extri_path, "w") as f:
        f.write("%YAML:1.0\n---\n")
        f.write(f"names:\n")
        for n in names:
            f.write(f'  - "{n}"\n')

        for cam_key, name in zip(cam_names, names):
            cam = vggt[cam_key]
            R_mat = np.array(cam["R"], dtype=np.float64)
            t_vec = np.array(cam["t"], dtype=np.float64)

            # Rodrigues vector from rotation matrix
            rvec = rotation_matrix_to_rodrigues(R_mat)

            write_opencv_matrix(f, f"R_{name}",
                                rvec.reshape(3, 1), 3, 1)
            write_opencv_matrix(f, f"Rot_{name}",
                                R_mat, 3, 3)
            write_opencv_matrix(f, f"T_{name}",
                                t_vec.reshape(3, 1), 3, 1)

    print(f"Written: {extri_path}")

    # ---- Verification ----
    print("\n=== Verification ===")
    for cam_key, name in zip(cam_names, names):
        cam = vggt[cam_key]
        K = np.array(cam["K"])
        R = np.array(cam["R"])
        t = np.array(cam["t"])

        # Camera center in world coordinates
        center = -R.T @ t
        rvec = rotation_matrix_to_rodrigues(R)

        print(f"cam{name}: fx={K[0,0]:.1f} fy={K[1,1]:.1f} "
              f"| center=[{center[0]:.3f},{center[1]:.3f},"
              f"{center[2]:.3f}] "
              f"| rvec=[{rvec[0]:.4f},{rvec[1]:.4f},"
              f"{rvec[2]:.4f}]")

    # Rodrigues roundtrip check
    R_orig = np.array(vggt["cam2"]["R"], dtype=np.float64)
    rvec = rotation_matrix_to_rodrigues(R_orig)
    R_back, _ = cv2.Rodrigues(rvec)
    err = np.max(np.abs(R_orig - R_back))
    print(f"\nRodrigues roundtrip error: {err:.2e}")


if __name__ == "__main__":
    main()
\end{lstlisting}


\section{실행 스크립트}

\begin{lstlisting}[caption={run\_vggt\_easymocap.sh --- 전체 파이프라인 실행}, label=lst:run, language=bash]
#!/bin/bash
set -e

# === Step 1: Convert VGGT -> EasyMocap format ===
echo "[Step 1] Converting VGGT calibration to EasyMocap format..."
cd /home/elicer
python3 vggt_to_easymocap.py \
  --input /home/elicer/vggt_calibration_result.json \
  --output_dir /home/elicer/easymocap_data

# Backup old calibration
cp /home/elicer/easymocap_data/intri.yml \
   /home/elicer/easymocap_data/intri_da3_backup.yml 2>/dev/null || true
cp /home/elicer/easymocap_data/extri.yml \
   /home/elicer/easymocap_data/extri_da3_backup.yml 2>/dev/null || true

echo "[Step 1] Done. intri.yml and extri.yml updated."

# === Step 2: Run EasyMocap triangulation + SMPL fitting ===
echo "[Step 2] Running EasyMocap..."
cd /home/elicer/EasyMocap

python3 apps/demo/mv1p.py \
  /home/elicer/easymocap_data \
  --out /home/elicer/easymocap_data/output/vggt_smpl \
  --cfg config/exp/mv1p/detect_triangulate_fitSMPL.yml \
  --vis_det --vis_repro

echo "[Step 2] Done. Results in output/vggt_smpl/"

# === Step 3: Quick validation ===
echo "[Step 3] Checking SMPL output..."
ls -la /home/elicer/easymocap_data/output/vggt_smpl/smpl/
head -50 /home/elicer/easymocap_data/output/vggt_smpl/smpl/000000.json

echo "=== Pipeline complete ==="
\end{lstlisting}


\section{검증 스크립트: 재투영 오차}

\begin{lstlisting}[caption={validate\_reprojection.py --- 삼각측량 결과 검증}, label=lst:validate]
#!/usr/bin/env python3
"""
Validate triangulation quality by computing reprojection error.
"""

import json
import numpy as np
from pathlib import Path


def load_cameras_yml(intri_path, extri_path):
    """Load EasyMocap camera YAML files."""
    import cv2
    fs_intri = cv2.FileStorage(str(intri_path),
                               cv2.FILE_STORAGE_READ)
    fs_extri = cv2.FileStorage(str(extri_path),
                               cv2.FILE_STORAGE_READ)

    cameras = {}
    for i in range(1, 8):
        name = str(i)
        K = fs_intri.getNode(f"K_{name}").mat()
        dist = fs_intri.getNode(f"dist_{name}").mat()
        Rot = fs_extri.getNode(f"Rot_{name}").mat()
        T = fs_extri.getNode(f"T_{name}").mat()

        cameras[name] = {
            "K": K, "dist": dist,
            "R": Rot, "t": T.flatten()
        }

    fs_intri.release()
    fs_extri.release()
    return cameras


def project_3d_to_2d(point3d, K, R, t):
    """Project a 3D point to 2D pixel coordinates."""
    p_cam = R @ point3d + t
    p_proj = K @ p_cam
    return p_proj[:2] / p_proj[2]


def compute_reprojection_error(cameras, kp3d, kp2d_per_cam):
    """Compute mean reprojection error across all cameras."""
    errors = []
    for cam_name, cam in cameras.items():
        if cam_name not in kp2d_per_cam:
            continue
        kp2d = kp2d_per_cam[cam_name]  # (N, 3): u, v, conf
        K = cam["K"]
        R = cam["R"]
        t = cam["t"]

        for j in range(len(kp3d)):
            if kp2d[j][2] < 0.3:  # low confidence
                continue
            proj = project_3d_to_2d(kp3d[j], K, R, t)
            err = np.linalg.norm(proj - kp2d[j][:2])
            errors.append(err)

    return np.mean(errors), np.median(errors), errors
\end{lstlisting}


% ====================================================================
\chapter{개발 방향 로드맵}
% ====================================================================

\begin{table}[H]
\centering
\caption{개발 로드맵}
\renewcommand{\arraystretch}{1.5}
\begin{tabular}{c p{7cm} p{4cm}}
\toprule
\textbf{Phase} & \textbf{작업} & \textbf{성공 기준} \\
\midrule
1 & VGGT $\rightarrow$ intri.yml / extri.yml 변환 & 포맷 정상, roundtrip 검증 \\
2 & EasyMocap 삼각측량 실행 & reprojection error $< 10$ px \\
3 & EasyMocap SMPL 피팅 실행 & MPJPE $< 50$ mm \\
4 & SMPL 메시 오버레이 시각화 & 시각적으로 인체에 정렬됨 \\
5 & 스케일 보정 (필요시) & 신장/사지 비율 현실적 \\
6 & GART 입력 준비 ($\vtheta, \vbeta$ per frame) & GART 학습 시작 가능 \\
\midrule
\multicolumn{3}{l}{\textbf{Fallback}: Phase 2--3 실패 시} \\
\midrule
F1 & VGGT K로 직접 삼각측량 (EasyMocap 우회) & 3D 키포인트 시각적 검증 \\
F2 & smplx 직접 피팅 (PyTorch, 서버 GPU) & MPJPE $< 50$ mm \\
F3 & 3DGS-Avatar JMBO (HMR 초기값 + joint opt.) & photometric loss 수렴 \\
\bottomrule
\end{tabular}
\end{table}


\section{잠재적 문제와 대비책}

\begin{enumerate}[leftmargin=*]
  \item \textbf{VGGT 스케일 문제}: VGGT의 translation이 정규화되어 있어 EasyMocap의 3D 재구성 스케일이 현실과 다를 수 있음.
    \begin{itemize}
      \item 대비: 피사체 키 또는 어깨 너비로 스케일 보정
      \item SMPL 피팅 자체는 상대적 관절 각도이므로 스케일 영향 적음
    \end{itemize}

  \item \textbf{EasyMocap의 SMPL 모델}: 기존에 \texttt{SMPL\_NEUTRAL.pkl} (304KB 경량)을 사용하여 뼈 길이가 부정확했음.
    \begin{itemize}
      \item 대비: \texttt{SMPL\_NEUTRAL.pkl} (full, $\sim$21MB)로 교체
    \end{itemize}

  \item \textbf{투구 동작의 극단적 포즈}: EasyMocap의 기본 pose prior가 일상 동작 기준이라 7,000 deg/s 어깨 회전에 적합하지 않을 수 있음.
    \begin{itemize}
      \item 대비: \texttt{lambda\_theta}를 0.0001로 낮추어 극단 포즈 허용
    \end{itemize}

  \item \textbf{왜곡 미보정}: VGGT는 왜곡 계수를 제공하지 않으므로, 이미지 가장자리의 키포인트에서 수 px 오차 발생 가능.
    \begin{itemize}
      \item 대비: 투수가 이미지 중앙에 있으면 영향 최소 ($< 5$ px)
      \item 필요시 OpenCV \texttt{calibrateCamera}로 사후 왜곡 추정 가능
    \end{itemize}
\end{enumerate}


% ====================================================================
\chapter{기대 효과}
% ====================================================================

\begin{summary}
이전 실패 (DA3 K $\rightarrow$ EasyMocap)와 비교한 개선점:

\begin{enumerate}[leftmargin=*]
  \item \textbf{좌표계 일치}: VGGT K가 원본 해상도 $\Leftrightarrow$ OpenPose 2D 키포인트 --- 역변환 불필요, 불일치 위험 제거
  \item \textbf{$f_x/f_y$ 정확도}: DA3 최대 2.2\% 편차 $\rightarrow$ VGGT 최대 0.6\% 편차
  \item \textbf{다시점 일관성}: VGGT가 7장을 동시에 보고 R, t 추정 $\rightarrow$ 카메라 간 상대 위치가 더 일관적
  \item \textbf{Y/Z 축 교환 해소}: DA3의 PnP 좌표계와 EasyMocap 좌표계 불일치 $\rightarrow$ VGGT는 OpenCV 규약으로 통일
\end{enumerate}

이전 MPJPE 219.6 mm에서 \textbf{50 mm 이하}를 목표로 한다.
\end{summary}


\end{document}
